\documentclass{article}
\usepackage{amsmath, amssymb, amsfonts, bm}
\usepackage{geometry}
\geometry{margin=1in}
\title{Teleparallel Cube Geometry and Prime-Torsion Recursion}
\author{Multiplicity Theory Project}
\date{June 2025}

\begin{document}
\maketitle

\section*{1. Teleparallel Cube Manifold Setup}

Let spacetime be discretized as a teleparallel unit cube centered at the origin. Define body diagonals:
\begin{align*}
\vec{d}_0 &= (1,1,1) &\text{(time-like)} \\
\vec{d}_1 &= (1,-1,1), &\vec{d}_2 = (-1,1,1), &\vec{d}_3 = (1,1,-1) &\text{(space-like)}
\end{align*}

We introduce a tetrad field $e^a{}_{\mu}$ that aligns with these directions:
\begin{equation}
    e^a{}_{\mu} = \frac{1}{\sqrt{3}}
    \begin{bmatrix}
        1 & 1 & 1 & 0 \\
        1 & -1 & 1 & 0 \\
        -1 & 1 & 1 & 0 \\
        1 & 1 & -1 & 0
    \end{bmatrix}
\end{equation}
Each row corresponds to a distinct diagonal vector in 3-space, extended trivially in time.

\subsection*{1.1 Metric from Tetrads}

Given $\eta_{ab} = \text{diag}(-1,+1,+1,+1)$, compute:
\begin{equation}
    g_{\mu\nu} = \eta_{ab} e^a{}_{\mu} e^b{}_{\nu}
\end{equation}
This yields a symmetric metric representing a locally flat but torsion-rich geometry.

\section*{2. Torsion Tensor and Scalar}

The Weitzenb\"ock connection is defined as:
\begin{equation}
    \Gamma^\lambda{}_{\mu\nu} = e_a{}^\lambda (\partial_\nu e^a{}_{\mu})
\end{equation}
From this we define the torsion tensor:
\begin{equation}
    T^\lambda{}_{\mu\nu} = \Gamma^\lambda{}_{\nu\mu} - \Gamma^\lambda{}_{\mu\nu}
\end{equation}

\subsection*{2.1 Torsion Scalar}
We define the torsion scalar:
\begin{equation}
    T = \frac{1}{4} T^\lambda{}_{\mu\nu} T_\lambda{}^{\mu\nu} + \frac{1}{2} T^\lambda{}_{\mu\nu} T^{\nu\mu}{}_{\lambda} - T^\rho{}_{\rho\lambda} T^{\nu\lambda}{}_{\nu}
\end{equation}

\subsection*{2.2 Superpotential and Torsion Force}
The superpotential is:
\begin{equation}
    S_\lambda^{\mu\nu} = \frac{1}{2} \left( K^{\mu\nu}{}_{\lambda} + \delta^\mu_\lambda T^{\alpha\nu}{}_{\alpha} - \delta^\nu_\lambda T^{\alpha\mu}{}_{\alpha} \right)
\end{equation}
with contortion tensor $K^{\mu\nu}{}_{\lambda} = -\frac{1}{2}(T^{\mu\nu}{}_{\lambda} - T^{\nu\mu}{}_{\lambda} - T_\lambda{}^{\mu\nu})$.

The force 4-vector becomes:
\begin{equation}
    F^\mu = \tav \cdot \partial_\nu T^\lambda{}_{\nu\lambda}
\end{equation}
This links torsional divergence to physical force effects modulated by $\tav = \hbar / c$.

\section*{3. Prime-Indexed Recursive Operator and Generator}

Let $p_i$ be the $i$-th prime. Define recursive evolution:
\begin{equation}
    \Xi^\mu(t+1; p_i) = \mathcal{F}_{p_i}[\Xi^\mu(t)]
\end{equation}
with:
\begin{equation}
    \mathcal{F}_{p_i}[\Xi] = \exp\left(\alpha_{p_i} L_{T^a} \right) \Xi, \quad L_{T^a} = T^a{}^\mu \partial_\mu
\end{equation}

Here, $L_{T^a}$ is a Lie-type generator along the torsion field direction, and $\alpha_{p_i}$ modulates its strength via the prime index. For Hecke-like interpretation:
\begin{equation}
    \mathcal{F}_{p_i}[\Xi] = T(p_i) \cdot \Xi, \quad T(p_i) \text{ acting as a Hecke operator on a modular torsion lattice.}
\end{equation}

\section*{Next Steps}
- Implement symbolic forms in Mathematica:
  - Define $e^a{}_{\mu}$ and compute $T^\lambda{}_{\mu\nu}$.
  - Evaluate $\mathcal{F}_{p_i}$ as both Lie and Hecke generators.
  - Discretize $\Xi(t; p_i)$ and evolve on a lattice.
\nocite{*}
\bibliographystyle{unsrt}
\bibliography{references}
\end{document}
